
\documentclass[addpoints]{exam}
\usepackage[version=4]{mhchem}
\usepackage{siunitx}
\usepackage{color}
\usepackage{amsmath}

% \printanswers

\newcommand{\event}{Structure and Properties of Matter}
\newcommand{\tournament}{}  % Replace with class name, date, or course code
\def\type{0}

\setlength\fillinlinelength{\textwidth}
\CorrectChoiceEmphasis{\color{red}\bfseries}
\SolutionEmphasis{\color{red}\bfseries}

\setlength\answerclearance{2pt}
\pagestyle{headandfoot}
\runningheadrule
\runningheader{\event}{\tournament}{\ifnum\type=1 Class Set Test \else Full Name:\hspace{1cm} \fi}
\runningfootrule
\runningfooter{}{Page \thepage\ of \numpages}{}

\begin{document}
\begin{center}
    {\huge Grade 12 Chemistry \\ \event
    \ifprintanswers
     \space Key
    \else
     \space \ifnum\type<2 Test \fi \ifnum\type=1 \textbf{(CLASS SET)} \fi \ifnum\type=2 Answer Sheet \fi
    \fi \\ \vspace{3mm}\tournament\vspace{1mm}}
\end{center}

\vspace{.25\textheight}

\begin{center}
    First Name: \ifprintanswers
    \underline{\hspace{3.5cm}}\textbf{KEY}\underline{\hspace{5.5cm}}\\\vspace{5mm}
    \else
        \underline{\hspace{10cm}}\\ \vspace{5mm}
    \fi
    Last Name: \underline{\hspace{10cm}}\\ \vspace{5mm}
\end{center}

\noindent\textbf{\underline{Directions:}}
\begin{itemize}
    \ifnum\type=1 \item \textbf{This test is a class set.} Please do not write any answers on this paper. \fi
    \item Show all work for full marks on short answer questions.
    \item Express answers to the appropriate number of significant figures.
    \item The test is designed to be completed in 75 minutes.
\end{itemize}
\begin{center}
\textit{For grading use only}\\
\multirowgradetable{1}[pages]
\end{center}

\newpage
\begin{questions}

\section*{Part A --- Multiple Choice (15 marks)}
\vspace{5mm}

\question[1] Which set of quantum numbers is valid for an electron?
\begin{choices}
    \choice $n=2,\ l=2,\ m_l=0,\ m_s=+\tfrac{1}{2}$
    \choice $n=3,\ l=1,\ m_l=-2,\ m_s=-\tfrac{1}{2}$
    \correctchoice $n=3,\ l=2,\ m_l=1,\ m_s=+\tfrac{1}{2}$
    \choice $n=1,\ l=1,\ m_l=0,\ m_s=-\tfrac{1}{2}$
\end{choices}

\question[1] How many electrons can occupy the 3d subshell?
\begin{choices}
    \choice 2
    \choice 6
    \choice 8
    \correctchoice 10
\end{choices}

\question[1] What is the correct ground-state electron configuration of \ce{Fe^{2+}}?
(\ce{Fe}: $[\text{Ar}]3d^6 4s^2$)
\begin{choices}
    \choice $[\text{Ar}]3d^4 4s^2$
    \choice $[\text{Ar}]3d^5 4s^1$
    \correctchoice $[\text{Ar}]3d^6$
    \choice $[\text{Ar}]3d^6 4s^2$
\end{choices}

\question[1] Which of the following elements has the highest first ionization energy?
\begin{choices}
    \choice Na
    \choice Mg
    \choice Al
    \correctchoice Ne
\end{choices}

\question[1] Which of the following atoms has the smallest atomic radius?
\begin{choices}
    \choice Na
    \choice Mg
    \choice Al
    \correctchoice Cl
\end{choices}

\question[1] What is the molecular geometry of \ce{SF4}?
\begin{choices}
    \choice Tetrahedral
    \correctchoice See-saw
    \choice Square planar
    \choice Trigonal pyramidal
\end{choices}

\question[1] Which of the following bonds is the most polar?
\begin{choices}
    \choice \ce{C-H}
    \choice \ce{N-H}
    \choice \ce{O-H}
    \correctchoice \ce{F-H}
\end{choices}

\question[1] What is the hybridization of each carbon atom in ethene (\ce{C2H4})?
\begin{choices}
    \choice $sp$
    \correctchoice $sp^2$
    \choice $sp^3$
    \choice $sp^3d$
\end{choices}

\question[1] Which type of intermolecular force is present in \textit{all} molecular substances?
\begin{choices}
    \choice Hydrogen bonding
    \choice Dipole--dipole forces
    \correctchoice London dispersion forces
    \choice Ion--dipole forces
\end{choices}

\question[1] Which of the following molecules is capable of forming hydrogen bonds?
\begin{choices}
    \choice \ce{CH4}
    \choice \ce{PH3}
    \choice \ce{H2S}
    \correctchoice \ce{HF}
\end{choices}

\question[1] Which of the following compounds has the highest boiling point primarily due to London dispersion forces?
\begin{choices}
    \choice \ce{CH4}
    \choice \ce{C2H6}
    \choice \ce{C3H8}
    \correctchoice \ce{C4H10}
\end{choices}

\question[1] What is the formal charge on nitrogen in \ce{NH4+}?
\begin{choices}
    \choice $-1$
    \choice $0$
    \correctchoice $+1$
    \choice $+2$
\end{choices}

\question[1] Which statement about resonance structures is correct?
\begin{choices}
    \choice Each resonance structure exists separately at different times.
    \choice The molecule alternates rapidly between resonance structures.
    \correctchoice The actual molecule is a hybrid of all resonance structures.
    \choice Resonance structures always have different bond lengths.
\end{choices}

\question[1] Which type of solid generally has the highest melting point?
\begin{choices}
    \choice Molecular solid
    \choice Metallic solid
    \choice Ionic solid
    \correctchoice Network covalent solid
\end{choices}

\question[1] Which of the following molecules is nonpolar despite containing polar bonds?
\begin{choices}
    \choice \ce{H2O}
    \choice \ce{NH3}
    \correctchoice \ce{CO2}
    \choice \ce{HCl}
\end{choices}


\section*{Part B --- Short Answer (33 marks)}

\question Chromium (Cr, $Z = 24$) has an anomalous electron configuration.
\begin{parts}
    \part[3] Write the full electron configuration of Cr and explain why it differs from the
    expected $[\text{Ar}]3d^4 4s^2$ pattern.
    \part[2] Write the electron configuration of \ce{Cr^{3+}} and state the number of unpaired
    electrons.
    \part[3] Compare the first ionization energies of Cr, Mn, and Fe across Period 4.
    Predict which has the highest value and explain the trend, noting any anomalies.
\end{parts}
\begin{solution}[9cm]
\textbf{(a)} Actual: $[\text{Ar}]3d^5 4s^1$. A half-filled $3d^5$ subshell has extra stability due
to maximized exchange energy and symmetric electron distribution. Promoting one $4s$ electron to
$3d$ achieves this stability, lowering the overall energy of the atom.

\textbf{(b)} \ce{Cr^{3+}} loses 3 electrons (4s first, then 3d): $[\text{Ar}]3d^3$.
Number of unpaired electrons: \textbf{3}.

\textbf{(c)} Across Period 4, first ionization energy generally increases. Mn ($3d^5 4s^2$) has a
half-filled $3d$ subshell with extra stability, so its IE is slightly higher than expected --- higher
than Fe. Approximate order: Fe $<$ Cr $<$ Mn. Mn has the highest IE among the three.
\end{solution}

\question Consider the molecule \ce{PCl5}.
\begin{parts}
    \part[2] Determine the electron geometry and molecular geometry of \ce{PCl5}.
    \part[2] Identify the hybridization of the phosphorus atom and explain why phosphorus can
    expand its octet but nitrogen cannot.
    \part[2] Is \ce{PCl5} polar or nonpolar? Justify your answer using bond dipoles.
    \part[2] \ce{PCl3} is also a known compound. Compare its molecular geometry and polarity
    to \ce{PCl5}.
\end{parts}
\begin{solution}[9cm]
\textbf{(a)} P has 5 bonding pairs and 0 lone pairs.
Electron geometry: \textbf{trigonal bipyramidal}.
Molecular geometry: \textbf{trigonal bipyramidal}.

\textbf{(b)} $sp^3d$ hybridization. Phosphorus can expand its octet because it is in Period 3
and has available, low-energy 3d orbitals. Nitrogen is in Period 2 and has no accessible $d$
orbitals, so it is limited to an octet.

\textbf{(c)} \textbf{Nonpolar.} The five P--Cl bond dipoles cancel: three equatorial bonds are
symmetric at $120°$ and two axial bonds are opposite each other at $180°$. The net dipole moment
is zero.

\textbf{(d)} \ce{PCl3}: P has 3 bonding pairs + 1 lone pair. Electron geometry: tetrahedral;
molecular geometry: \textbf{trigonal pyramidal}. The lone pair creates an asymmetric charge
distribution, making \ce{PCl3} \textbf{polar}. Unlike \ce{PCl5}, its bond dipoles do not cancel.
\end{solution}

\question Compare the intermolecular forces and boiling points of the following substances.
\begin{parts}
    \part[4] Identify all types of intermolecular forces present in each of the following
    molecules: \ce{CH4},\ \ce{NH3},\ \ce{H2O},\ \ce{HF}.
    \part[3] Explain why \ce{H2O} has a much higher boiling point (\SI{100}{\celsius}) than
    \ce{H2S} (\SI{-61}{\celsius}), despite \ce{H2O} having a significantly lower molar mass.
    \part[2] Arrange Ne, Ar, Kr, Xe in order of increasing boiling point and justify your
    answer.
\end{parts}
\begin{solution}[9cm]
\textbf{(a)}
\begin{itemize}
    \item \ce{CH4}: London dispersion forces only (nonpolar, no H bonded to F, O, or N).
    \item \ce{NH3}: London dispersion, dipole--dipole, and \textbf{hydrogen bonding}
    (N--H$\cdots$N).
    \item \ce{H2O}: London dispersion, dipole--dipole, and \textbf{hydrogen bonding}
    (O--H$\cdots$O); strongest network because O has high electronegativity and two O--H bonds per
    molecule.
    \item \ce{HF}: London dispersion, dipole--dipole, and \textbf{hydrogen bonding}
    (F--H$\cdots$F).
\end{itemize}

\textbf{(b)} Each \ce{H2O} molecule can donate 2 hydrogen bonds (two O--H bonds) and accept 2
(two lone pairs on O), forming an extensive 3-D hydrogen-bonding network. Oxygen's high
electronegativity ($\chi = 3.44$) makes O--H bonds very polar. Sulfur in \ce{H2S} has lower
electronegativity and a larger atomic radius; S--H bonds are too weak and too long to form true
hydrogen bonds. Only London dispersion and weak dipole--dipole forces act between \ce{H2S}
molecules.

\textbf{(c)} Ne $<$ Ar $<$ Kr $<$ Xe. Larger noble gas atoms have more electrons and greater
polarizability, leading to stronger London dispersion forces and higher boiling points.
\end{solution}

\question Consider the following four species: \ce{BeCl2}, \ce{SO3}, \ce{XeF4}, \ce{IF5}.
\begin{parts}
    \part[4] For each species, state the (i) electron geometry, (ii) molecular geometry, and
    (iii) whether it is polar or nonpolar.
    \part[4] Draw all three resonance structures of \ce{SO3} and calculate the bond order of
    each \ce{S-O} bond. Explain what the bond order tells you about bond length and strength
    compared to a pure \ce{S-O} single bond or \ce{S=O} double bond.
\end{parts}
\begin{solution}[10cm]
\textbf{(a)}
\begin{itemize}
    \item \ce{BeCl2}: 2 bonding pairs, 0 lone pairs on Be. Electron geo: \textbf{linear};
    Molecular geo: \textbf{linear}; \textbf{Nonpolar} (symmetric, dipoles cancel).
    \item \ce{SO3}: 3 bonding regions, 0 lone pairs on S. Electron geo: \textbf{trigonal
    planar}; Molecular geo: \textbf{trigonal planar}; \textbf{Nonpolar} (symmetric arrangement).
    \item \ce{XeF4}: 4 bonding pairs + 2 lone pairs on Xe. Electron geo: \textbf{octahedral};
    Molecular geo: \textbf{square planar}; \textbf{Nonpolar} (lone pairs are trans to each other,
    bond dipoles cancel).
    \item \ce{IF5}: 5 bonding pairs + 1 lone pair on I. Electron geo: \textbf{octahedral};
    Molecular geo: \textbf{square pyramidal}; \textbf{Polar} (lone pair breaks symmetry).
\end{itemize}

\textbf{(b)} \ce{SO3} has 3 equivalent resonance structures, each showing one \ce{S=O} double
bond and two \ce{S-O} single bonds in different positions.

\[
\text{Bond order} = \frac{\text{(number of bonds in all structures for one S--O position)}}
{\text{number of resonance structures}} = \frac{1\times2 + 1\times1 + 1\times1}{3} = \frac{4}{3}
\approx 1.33
\]

A bond order of $4/3$ means each \ce{S-O} bond is intermediate between a single bond (order 1)
and a double bond (order 2). All three bonds are equivalent: shorter and stronger than a pure
single bond, longer and weaker than a pure double bond.
\end{solution}

\end{questions}
\end{document}
